\documentclass[10pt,letterpaper]{article}

\usepackage{PRIMEarxiv}

\usepackage[utf8]{inputenc}
\usepackage[T1]{fontenc}
\usepackage{hyperref}
\usepackage{url}
\usepackage{booktabs}
\usepackage{amsfonts}
\usepackage[expansion=false]{microtype}
\usepackage{graphicx}
\usepackage{amsmath,amssymb}
\usepackage{tikz}
\usepackage{pgfplots}
\usetikzlibrary{shapes.geometric,arrows,positioning,calc}
\pgfplotsset{compat=1.18}

\usepackage{listings}
\lstset{%
  basicstyle=\ttfamily\small,
  breaklines=true,
  frame=single,
}

\def\figdir{figures}

\hypersetup{%
  pdftitle={Format-Aware Checkpoint Loading: An Adaptive Backend Policy for LLMs},
  pdfauthor={Botir Khaltaev},
  pdfkeywords={large language models, checkpoint loading, io\_uring, cold start, vLLM, reproducibility}%
}

%% Title
\title{Format-Aware Checkpoint Loading: An Adaptive Backend Policy for LLMs}

%% Authors
\author{
  Botir Khaltaev\\
  Independent, United Kingdom\\
  \texttt{btrghstk@gmail.com}
}

\begin{document}
\maketitle

\begin{abstract}
Inference on large language models requires loading weights from storage into GPU memory. In serverless settings where idle models are evicted, cold-start latency hinges on how checkpoints are read. Peak NVMe bandwidth is only an upper bound: checkpoint layout, request structure, and integration overhead often dominate realised throughput.

This paper treats checkpoint loading as a measurement-driven systems problem. Using Tensora, we compare synchronous POSIX I/O, Tokio async, io\_uring, and an adaptive \texttt{default} policy across contiguous SafeTensors and partitioned ServerlessLLM-style layouts, then trace those choices through Python bindings and vLLM integration.

For contiguous SafeTensors, explicit sync wins at smaller scales while explicit io\_uring wins at the largest multi-shard scales where queue depth matters. For partitioned ServerlessLLM, sync never leads; the relevant trade-off is among async, io\_uring, and \texttt{default}. These layout- and scale-dependent crossovers show that no single fixed backend is optimal in all cases.

Deploy a format-aware adaptive \texttt{default} with explicit overrides, not a single universal loader. Python and vLLM results confirm that integration overhead can reorder backend rankings; loader choice matters most for initialisation and time-to-first-token.
\end{abstract}

\keywords{large language models \and checkpoint loading \and I/O policy \and io\_uring \and cold start \and vLLM \and reproducibility}

\input{sections/introduction}
\input{sections/related_work}
\input{sections/design}
\input{sections/experimental_setup}
\input{sections/results}
\input{sections/discussion}
\input{sections/limitations}
\input{sections/professional_issues}
\input{sections/conclusion}

\appendix
\input{sections/appendix_reproducibility}

\bibliographystyle{unsrt}
\bibliography{library}

\end{document}
